\chapter{Topology of Metric Spaces}\label{ch:b2:metric}

The topology of the real line (Year 1 volume) generalizes, almost
without changing a word, to any set equipped with a distance. The
gain is enormous: sequences of functions, matrices, curves --- all
become points of \omterm{def:b2:metric:def}{metric spaces}, and the three pillars proved here
--- \omterm{def:b2:metric:complete}{completeness} with the Banach fixed point theorem, \omterm{def:b2:metric:compact}{compactness},
\omterm{def:b2:metric:connected}{connectedness} --- apply to them uniformly. This chapter is the
backbone of the whole analysis half of the book.

\section{Metric spaces}

\begin{definition}\label{def:b2:metric:def}
A \emph{metric space}\index{metric space} is a set $X$ with a map $d
\colon X \times X \to \R_+$ such that, for all $x, y, z$:
\[
d(x,y) = 0 \iff x = y,
\qquad
d(x,y) = d(y,x),
\qquad
d(x,z) \leq d(x,y) + d(y,z).
\]
Balls: $B(a, r) = \{x : d(a,x) < r\}$ (\omterm{def:b2:metric:topology}{open}), $\overline B(a,r) =
\{x : d(a,x) \leq r\}$ (closed). A subset $A \subseteq X$ becomes a
metric space with the induced distance.
\end{definition}

\begin{example}\label{ex:b2:metric:examples}
$\R$ with $\abs{x - y}$; $\R^n$ with any of
\[
d_1(x,y) = \sum_i \abs{x_i - y_i},
\quad
d_2(x,y) = \Bigl(\sum_i (x_i - y_i)^2\Bigr)^{1/2},
\quad
d_\infty(x,y) = \max_i \abs{x_i - y_i};
\]
the set $C(\intcc{a}{b})$ of \omterm{def:b2:metric:continuity}{continuous} functions with the
\emph{sup distance} $d_\infty(f, g) = \sup_{\intcc{a}{b}} \abs{f -
g}$ (finite: $f - g$ is bounded); any set with the \emph{discrete}
distance ($d(x,y) = 1$ for $x \neq y$). Distances coming from norms
are the subject of \cref{ch:b2:nvs}.
\end{example}

\begin{definition}[Topology of a metric space]\label{def:b2:metric:topology}
$U \subseteq X$ is \emph{open}\index{open set!in a metric space}
when every point of $U$ is the center of a ball contained in $U$;
$F$ is \emph{closed} when its complement is open. Neighborhoods,
interior, closure, density, boundary are defined exactly as on the
real line (Year 1 volume), with balls replacing intervals, and the
statements proved there --- unions/intersections of open sets,
characterizations of interior and closure, closure as smallest
closed superset --- carry over with the same proofs. Open balls are
open, closed balls are closed (triangle inequality).
\end{definition}

\begin{example}[Interior, closure, boundary on one set]\label{ex:b2:metric:intclosure}
In $\R$, let $A = \intoc{0}{1} \cup \{2\}$. Interior:
$\intoo{0}{1}$ --- around any $x \in \intoo01$ a small ball
stays in $A$; around $1$, every ball $\intoo{1-r}{1+r}$ leaks
out of $A$ on the right, so $1$ is not interior; and the
isolated $2$ is not interior either. Closure:
$\intcc{0}{1} \cup \{2\}$ (the point $0$ is a limit of $A$,
nothing else is added). Boundary (closure minus interior):
$\{0, 1, 2\}$. Note the asymmetries worth remembering: an
endpoint can belong to a set without being interior ($1$), can
be adherent without belonging ($0$), and an isolated point is
its own boundary ($2$). The same bookkeeping runs verbatim in
any \omterm{def:b2:metric:def}{metric space}, with balls in place of intervals.
\end{example}

\begin{definition}[Limits, continuity]\label{def:b2:metric:continuity}
$x_n \to x$ in $X$ when $d(x_n, x) \to 0$. A map $f \colon X \to Y$
between \omterm{def:b2:metric:def}{metric spaces} is \emph{continuous}\index{continuous} at $a$ when
\[
\forall \varepsilon > 0,\ \exists\delta > 0,\quad
d_X(x, a) \leq \delta \implies d_Y\bigl(f(x), f(a)\bigr) \leq
\varepsilon ;
\]
equivalently (same proof as on $\R$), $f(x_n) \to f(a)$ for every
sequence $x_n \to a$. $f$ is \emph{Lipschitz}\index{Lipschitz} with constant $k$ when
$d_Y(f(x), f(y)) \leq k\, d_X(x, y)$ always --- then uniformly
continuous, hence continuous.
\end{definition}

\begin{theorem}[Global characterization of continuity]\label{thm:b2:metric:globalcontinuity}
$f \colon X \to Y$ is \omterm{def:b2:metric:continuity}{continuous} (at every point) if and only if the
preimage of every \omterm{def:b2:metric:topology}{open} set is \omterm{def:b2:metric:topology}{open} --- iff the preimage of every
closed set is closed.
\end{theorem}

\begin{proof}
($\Rightarrow$) Let $V \subseteq Y$ be \omterm{def:b2:metric:topology}{open} and $a \in f^{-1}(V)$:
some ball $B(f(a), \varepsilon) \subseteq V$; \omterm{def:b2:metric:continuity}{continuity} at $a$
provides $\delta$ with $f(B(a, \delta)) \subseteq B(f(a),
\varepsilon)$, so $B(a, \delta) \subseteq f^{-1}(V)$.

($\Leftarrow$) Given $a$ and $\varepsilon$: $f^{-1}\bigl(B(f(a),
\varepsilon)\bigr)$ is \omterm{def:b2:metric:topology}{open} and contains $a$, hence contains a ball
$B(a, \delta)$: that is the definition of \omterm{def:b2:metric:continuity}{continuity} at $a$. Closed
sets: complements (\cref{prop:b2:structures:images}).
\end{proof}

\section{Complete spaces}

\begin{definition}\label{def:b2:metric:complete}
A sequence $(x_n)$ is \emph{Cauchy} when $\sup_{p, q \geq N} d(x_p,
x_q) \to 0$ as $N \to \infty$. A \omterm{def:b2:metric:def}{metric space} is
\emph{complete}\index{complete metric space} when every Cauchy
sequence converges. Convergent $\Rightarrow$ Cauchy always; closed
subsets of complete spaces are complete, and complete subsets of any
space are closed (same proofs as on $\R$: Year 1 volume).
\end{definition}

\begin{example}[Cauchy without a limit]\label{ex:b2:metric:cauchynolimit}
In $X = \Q$ with the usual distance, the decimal truncations of
$\sqrt2$,
\[
x_0 = 1,\quad x_1 = 1.4,\quad x_2 = 1.41,\quad x_3 = 1.414,
\quad\dots
\]
satisfy $\abs{x_p - x_q} \leq 10^{-\min(p,q)}$: Cauchy in $\Q$.
A limit in $\Q$ would also be the limit in $\R$, namely
$\sqrt2 \notin \Q$: no limit exists in $X$. Incompleteness is
the presence of such ``phantom limits''; \omterm{def:b2:metric:complete}{completeness} of $\R$
was engineered in the Year 1 volume precisely to give every
Cauchy sequence a home.
\end{example}

\begin{theorem}\label{thm:b2:metric:rncomplete}
$\R^n$ (any of the three distances of
\cref{ex:b2:metric:examples}) and $\bigl(C(\intcc{a}{b}),
d_\infty\bigr)$ are \omterm{def:b2:metric:complete}{complete}.
\end{theorem}

\begin{proof}
$\R^n$: a Cauchy sequence is Cauchy in each coordinate (each
$\abs{x_i - y_i} \leq d(x,y)$ for all three distances), so each
coordinate converges (\omterm{def:b2:metric:complete}{completeness} of $\R$, Year 1 volume), and
coordinatewise convergence implies convergence for $d_\infty$
(finitely many coordinates), hence for all three (the three
distances dominate each other within constant factors: $d_\infty
\leq d_2 \leq d_1 \leq n\,d_\infty$).

$C(\intcc{a}{b})$: let $(f_n)$ be $d_\infty$-Cauchy. For each $x$,
$(f_n(x))$ is Cauchy in $\R$ ($\abs{f_p(x) - f_q(x)} \leq
d_\infty(f_p, f_q)$): converges to some $f(x)$. Passing to the
limit in $\abs{f_p(x) - f_q(x)} \leq \varepsilon$ (valid for $p, q
\geq N_\varepsilon$, all $x$) as $q \to \infty$: $\abs{f_p(x) -
f(x)} \leq \varepsilon$ for all $x$, i.e.\ $d_\infty(f_p, f) \leq
\varepsilon$: uniform convergence. The limit is \omterm{def:b2:metric:continuity}{continuous}: given
$\varepsilon$, pick $p$ with $\sup\abs{f_p - f} \leq \varepsilon$,
then use \omterm{def:b2:metric:continuity}{continuity} of $f_p$ at $a$ and the three-term split
\[
\abs{f(x) - f(a)} \leq \abs{f(x) - f_p(x)} + \abs{f_p(x) - f_p(a)}
+ \abs{f_p(a) - f(a)} \leq 3\varepsilon
\]
for $x$ close to $a$. (This ``$3\varepsilon$ argument'' returns as
the uniform-limit theorem of \cref{ch:b2:funcseq}.)
\end{proof}

\begin{example}[Open and closed sets recognized by continuity]\label{ex:b2:metric:recognize}
The global characterization
(\cref{thm:b2:metric:globalcontinuity}) is the everyday tool for
topological bookkeeping. In $\R^2$: the set $\{(x, y) : x^2 + y^2
< 1,\ y > x^3\}$ is \omterm{def:b2:metric:topology}{open} --- it is $g^{-1}(\intoo{-\infty}{1})
\cap h^{-1}(\intoo{0}{+\infty})$ for the \omterm{def:b2:metric:continuity}{continuous} $g(x,y) =
x^2 + y^2$ and $h(x, y) = y - x^3$, an intersection of two \omterm{def:b2:metric:topology}{open}
preimages. In $\bigl(C(\intcc01), d_\infty\bigr)$: the set of
functions with $f(0) = f(1)$ and $\int_0^1 f = 0$ is closed ---
the preimage of $\{(0,0)\}$ under the \omterm{def:b2:metric:continuity}{continuous} map $f \mapsto
\bigl(f(0) - f(1),\ \int_0^1 f\bigr)$ into $\R^2$ (each
coordinate is $1$-Lipschitz, as in \cref{exo:b2:metric:3}). The
method never draws a picture: exhibit a \omterm{def:b2:metric:continuity}{continuous} map, read the
set as a preimage, quote the theorem.
\end{example}

\begin{example}[A closed set defined by infinitely many conditions]\label{ex:b2:metric:lipclosed}
In $\bigl(C(\intcc01), d_\infty\bigr)$, the set
\[
L = \{f : \abs{f(x) - f(y)} \leq \abs{x - y}
\ \text{for all } x, y\}
\]
of $1$-Lipschitz functions is closed, although it is cut out by
uncountably many conditions: for each fixed pair $(x, y)$, the
map $f \mapsto \abs{f(x) - f(y)} - \abs{x - y}$ is \omterm{def:b2:metric:continuity}{continuous}
(evaluations are $1$-Lipschitz), so each single condition
defines a closed set, and $L$ is the \emph{intersection} of
this family --- an arbitrary intersection of closed sets is
closed. The same template certifies closedness for monotone
functions, convex functions, functions bounded by a fixed $g$:
uniform limits inherit every property expressible as a family
of closed pointwise constraints. What uniform limits do
\emph{not} automatically inherit --- differentiability, for
one --- is exactly what \cref{ch:b2:funcseq} must labour for.
\end{example}

\begin{theorem}[Banach fixed point theorem]\label{thm:b2:metric:banach}
Let $X$ be a nonempty \omterm{def:b2:metric:complete}{complete metric space} and $f \colon X \to X$ a
\emph{contraction}: \omterm{def:b2:metric:continuity}{Lipschitz} with constant $k < 1$. Then $f$ has a
unique fixed point $\ell$, and every orbit $x_{n+1} = f(x_n)$
converges to $\ell$, with\index{Banach fixed point theorem}
\[
d(x_n, \ell) \leq \frac{k^n}{1 - k}\, d(x_1, x_0) .
\]
\end{theorem}

\begin{proof}
Uniqueness: two fixed points are at distance $\leq k$ times itself.
Existence: $d(x_{n+1}, x_n) \leq k^n d(x_1, x_0)$ by induction, so
for $q > p$,
\[
d(x_q, x_p) \leq \sum_{j=p}^{q-1} d(x_{j+1}, x_j)
\leq d(x_1, x_0) \sum_{j \geq p} k^j
= \frac{k^p}{1-k}\, d(x_1, x_0) \xrightarrow[p\to\infty]{} 0 :
\]
Cauchy, hence convergent to some $\ell$; \omterm{def:b2:metric:continuity}{continuity} of $f$ passes
$x_{n+1} = f(x_n)$ to the limit: $\ell = f(\ell)$. The error bound
is the displayed estimate with $q \to \infty$.
\end{proof}

\begin{example}[An integral equation]\label{ex:b2:metric:integralequation}
On $X = C(\intcc{0}{1})$ (\omterm{def:b2:metric:complete}{complete},
\cref{thm:b2:metric:rncomplete}), consider $T(f)(x) = 1 + \frac12
\int_0^x f(t)\,\dd t$. For $f, g \in X$:
\[
\abs{T(f)(x) - T(g)(x)} \leq \frac12 \int_0^x \abs{f - g} \leq
\frac12\, d_\infty(f, g),
\]
so $T$ is a $\frac12$-contraction: it has a unique \omterm{def:b2:metric:continuity}{continuous} fixed
point --- the solution of $f' = \frac f2$, $f(0) = 1$, namely
$\eu^{x/2}$. This scheme, industrialized, becomes the
Cauchy--Lipschitz theorem of \cref{ch:b2:diffeq}.
\end{example}

\begin{example}[A numerical fixed point: $x = \cos x$]\label{ex:b2:metric:cosfix}
On the \omterm{def:b2:metric:complete}{complete} $X = \intcc{0}{1}$, the map $f = \cos$ sends $X$
into $\intcc{\cos 1}{1} \subseteq X$ and is a contraction: by the
mean value inequality,
\[
\abs{\cos x - \cos y} \leq \bigl(\sup_{\intcc01}\abs{\sin}\bigr)
\abs{x - y} = (\sin 1)\abs{x - y},
\qquad \sin 1 \approx 0.841 < 1 .
\]
Banach: a unique solution of $x = \cos x$ in $\intcc01$ (hence in
$\R$: any real fixed point lies in $\intcc{-1}{1}$, then in
$\intcc{\cos 1}{1}$ after one application), and the iteration
$x_{n+1} = \cos x_n$ converges to it from any start: $x_\infty
\approx 0.739085$, the famous number obtained by hammering the
cosine key of a calculator. The error bound predicts
$(\sin1)^n/(1 - \sin1)$ decay --- about one digit per $13$
presses; the a posteriori bound of this chapter's weekend
problem (question 14) certifies each step on the fly.
\end{example}

\section{Compactness}

\begin{definition}\label{def:b2:metric:compact}
A \omterm{def:b2:metric:def}{metric space} $X$ is \emph{compact}\index{compact metric space}
when every sequence in $X$ has a subsequence converging \emph{in}
$X$ (the Bolzano--Weierstrass property). A subset is compact when it
is so with the induced distance.
\end{definition}

\begin{theorem}[First properties]\label{thm:b2:metric:compactprops}
\begin{enumerate}
  \item A \omterm{def:b2:metric:compact}{compact} subset is closed and bounded; a closed subset of a
        \omterm{def:b2:metric:compact}{compact} space is \omterm{def:b2:metric:compact}{compact}.
  \item In $\R^n$, the converse holds: \omterm{def:b2:metric:compact}{compact} $\iff$ closed and
        bounded.
  \item A \omterm{def:b2:metric:continuity}{continuous} image of a \omterm{def:b2:metric:compact}{compact} space is \omterm{def:b2:metric:compact}{compact}; a
        \omterm{def:b2:metric:continuity}{continuous} real function on a nonempty \omterm{def:b2:metric:compact}{compact} space is
        bounded and attains its bounds.
  \item (Heine) A \omterm{def:b2:metric:continuity}{continuous} map on a \omterm{def:b2:metric:compact}{compact} space is uniformly
        continuous.
  \item Products: if $X, Y$ are \omterm{def:b2:metric:compact}{compact}, so is $X \times Y$ (with
        $d\bigl((x,y),(x',y')\bigr) = d(x,x') + d(y,y')$).
\end{enumerate}
\end{theorem}

\begin{proof}
(1) Same arguments as on the line (Year 1 volume): an escaping-to-%
infinity or converging-outside sequence has no subsequence
converging inside; for the second claim, extract in the ambient
\omterm{def:b2:metric:compact}{compact} and use closedness.

(2) Bounded sequences in $\R^n$ have componentwise convergent
subsequences: extract on the first coordinate
(Bolzano--Weierstrass on $\R$), then, from that subsequence, on the
second, and so on ($n$ successive extractions); closedness keeps
the limit inside.

(3) Given $(f(x_n))$, extract $x_{\varphi(n)} \to x \in X$;
\omterm{def:b2:metric:continuity}{continuity} gives $f(x_{\varphi(n)}) \to f(x) \in f(X)$. Real case:
\omterm{def:b2:metric:compact}{compactness} of $f(X) \subseteq \R$ makes it closed and bounded, and
$\sup f(X) \in f(X)$ (the sup of a set is adherent to it, and
$f(X)$ is closed).

(4) The Year 1 proof transfers verbatim; here it is, in metric
dress. Suppose $f \colon X \to Y$ \omterm{def:b2:metric:continuity}{continuous} on the \omterm{def:b2:metric:compact}{compact} $X$
but not uniformly continuous: some $\varepsilon > 0$ admits, for
every $n$, points with
\[
d_X(x_n, y_n) \leq \frac{1}{n+1}
\qquad\text{and}\qquad
d_Y\bigl(f(x_n), f(y_n)\bigr) > \varepsilon .
\]
Extract $x_{\varphi(n)} \to a \in X$; then $y_{\varphi(n)} \to
a$ too (the mutual distances tend to $0$). \omterm{def:b2:metric:continuity}{Continuity} at $a$
sends both image sequences to $f(a)$, so
$d_Y\bigl(f(x_{\varphi(n)}), f(y_{\varphi(n)})\bigr) \to 0$ ---
contradicting the uniform gap $> \varepsilon$. \omterm{def:b2:metric:compact}{Compactness}
supplied exactly one thing: the cluster point $a$ at which to
apply plain \omterm{def:b2:metric:continuity}{continuity}.

(5) Extract on the $X$-coordinates, then again on the
$Y$-coordinates.
\end{proof}

\begin{example}[Heine's theorem, with and without compactness]\label{ex:b2:metric:heinework}
On $\intcc{0}{1}$, the function $x \mapsto x^2$ is uniformly
continuous --- Heine says so with no computation, but the direct
estimate is instructive:
\[
\abs{x^2 - y^2} = \abs{x + y}\,\abs{x - y} \leq 2\abs{x - y},
\]
so $\delta = \varepsilon/2$ works \emph{for every point at
once}. On $\R$ the same function is not uniformly continuous:
with $x_n = n$ and $y_n = n + \frac1n$, the gap $\abs{x_n - y_n}
= \frac1n \to 0$ while $\abs{x_n^2 - y_n^2} = 2 + \frac1{n^2}
\geq 2$: no single $\delta$ serves $\varepsilon = 1$. The
mechanism is visible: the local \omterm{def:b2:metric:continuity}{Lipschitz} constant $\abs{x + y}$
is bounded on a \omterm{def:b2:metric:compact}{compact} and unbounded on $\R$ --- Heine's
theorem is exactly the statement that \omterm{def:b2:metric:compact}{compactness} caps such
local constants uniformly.
\end{example}

\begin{method}[Proving that a set is compact]\label{met:b2:metric:compactness}
Three routes, in order of frequency. (1) \emph{Ambient
recognition:} in $\R^n$ (or any finite-dimensional normed
space, \cref{ch:b2:nvs}), verify closed --- typically as a
preimage, \cref{ex:b2:metric:recognize} --- and bounded. (2)
\emph{Inheritance:} a closed subset of a known \omterm{def:b2:metric:compact}{compact} space is
\omterm{def:b2:metric:compact}{compact}; a finite union or a product of \omterm{def:b2:metric:compact}{compacts} is \omterm{def:b2:metric:compact}{compact}; a
\omterm{def:b2:metric:continuity}{continuous} image of a \omterm{def:b2:metric:compact}{compact} is \omterm{def:b2:metric:compact}{compact}. (3) \emph{Bare
hands:} extract a convergent subsequence from an arbitrary
sequence --- usually by successive extractions coordinate by
coordinate. To prove \emph{non}-compactness, one witness
suffices: a sequence with no convergent subsequence, most often
points at mutual distance $\geq \varepsilon$.
\end{method}

\begin{example}[Distances between sets: compactness earns its keep]\label{ex:b2:metric:setdistance}
Let $K$ be \omterm{def:b2:metric:compact}{compact}, $F$ closed, $K \cap F = \emptyset$ in a
\omterm{def:b2:metric:def}{metric space}. Then
\[
d(K, F) = \inf\,\{d(x, y) : x \in K,\ y \in F\} > 0 :
\]
the function $x \mapsto d(x, F)$ is \omterm{def:b2:metric:continuity}{continuous}
(\cref{exo:b2:metric:11}) and positive on $K$ ($d(x, F) = 0$
would put $x \in \overline F = F$), so it attains a positive
minimum on the \omterm{def:b2:metric:compact}{compact} $K$
(\cref{thm:b2:metric:compactprops}~(3)). \omterm{def:b2:metric:compact}{Compactness} is not
decorative: for two \emph{closed} sets the infimum can vanish
without being attained --- in $\R^2$, the hyperbola $F_1 = \{xy
= 1\}$ and the axis $F_2 = \{y = 0\}$ are disjoint closed sets
with $d(F_1, F_2) = 0$ (the points $(n, \frac1n)$ approach the
axis). Escape to infinity is exactly what \omterm{def:b2:metric:compact}{compactness} forbids.
\end{example}

\begin{theorem}[Borel--Lebesgue]\label{thm:b2:metric:borellebesgue}
A \omterm{def:b2:metric:def}{metric space} $X$ is \omterm{def:b2:metric:compact}{compact} if and only if every cover of $X$ by
\omterm{def:b2:metric:topology}{open} sets has a \emph{finite} subcover.\index{Borel--Lebesgue
theorem}
\end{theorem}

\begin{proof}
($\Leftarrow$) Suppose $(x_n)$ has no convergent subsequence. We
claim every $x \in X$ has a ball $B(x, r_x)$ containing $x_n$
for only finitely many indices $n$: otherwise, every ball
$B(x, \frac1{k+1})$ would contain infinitely many terms, and
choosing indices
\[
\varphi(0) < \varphi(1) < \varphi(2) < \cdots
\quad\text{with}\quad
x_{\varphi(k)} \in B\Bigl(x, \frac{1}{k+1}\Bigr)
\]
(possible at each step precisely because infinitely many
candidates remain) would build a subsequence converging to $x$.
The balls $B(x, r_x)$ cover $X$; if finitely many of them
covered $X$, the index set $\N$ would be a finite union of
finite sets: absurd.

($\Rightarrow$) Two steps. \emph{Lebesgue number:} for an \omterm{def:b2:metric:topology}{open} cover
$(U_i)$ of a \omterm{def:b2:metric:compact}{compact} $X$, there is $\rho > 0$ such that every ball
of radius $\rho$ lies in some $U_i$. Otherwise, for each $n$ pick
$x_n$ with $B(x_n, \frac{1}{n+1})$ in no $U_i$; extract
$x_{\varphi(n)} \to x \in U_{i_0} \supseteq B(x, r)$; for large $n$,
$B(x_{\varphi(n)}, \frac{1}{\varphi(n)+1}) \subseteq B(x, r)
\subseteq U_{i_0}$: contradiction. \emph{Total boundedness:} for
every $\rho > 0$, finitely many balls of radius $\rho$ cover $X$.
Otherwise pick inductively $x_{n+1}$ outside $B(x_0, \rho) \cup
\dots \cup B(x_n, \rho)$: the sequence has pairwise distances $\geq
\rho$, so no Cauchy --- hence no convergent --- subsequence:
contradiction. Combining: cover $X$ by finitely many balls of radius
$\rho$ (the Lebesgue number), each inside some $U_i$: a finite
subcover.
\end{proof}

\begin{example}[An $\varepsilon$-net, counted]\label{ex:b2:metric:epsnet}
Total boundedness (from the proof of
\cref{thm:b2:metric:borellebesgue}) is very concrete on
$\intcc{0}{1}$: for $\varepsilon > 0$, the $\lceil
\frac{1}{2\varepsilon}\rceil$ balls centered at $\varepsilon,
3\varepsilon, 5\varepsilon, \dots$ of radius $\varepsilon$
cover it --- about $\frac1{2\varepsilon}$ balls, and no cover
can do with fewer than $\frac{1}{2\varepsilon}$ of them (each
ball covers length at most $2\varepsilon$). In $\intcc01^2$
the count squares to order $\varepsilon^{-2}$: covering numbers
grow like $\varepsilon^{-d}$ in dimension $d$ --- a quantitative
face of \omterm{def:b2:metric:compact}{compactness}, and the reason the infinite-dimensional
unit balls of \cref{ch:b2:nvs} (where no finite
$\frac13$-net exists at all) cannot be \omterm{def:b2:metric:compact}{compact}.
\end{example}

\begin{example}[Reading compactness on covers]\label{ex:b2:metric:covers}
The \omterm{def:b2:metric:topology}{half-open} interval $\intoc{0}{1}$ is covered by the \omterm{def:b2:metric:topology}{open} sets
$U_n = \intoo{\frac1n}{2}$, $n \geq 1$; any finite subfamily has a
largest index $N$ and misses $\intoc{0}{\frac1N}$: no finite
subcover, so $\intoc{0}{1}$ is not \omterm{def:b2:metric:compact}{compact} --- which the
sequential definition sees through $x_n = \frac1n$, whose limit
$0$ escapes. On the other hand, adding the single point $0$
repairs both diagnoses at once: on $\intcc{0}{1}$ every such cover
must contain a set containing $0$, which swallows a whole initial
segment, and finitely many sets finish the rest. The two languages
of \cref{thm:b2:metric:borellebesgue} always fail or succeed
together --- covers detect escape exactly where sequences do.
\end{example}

\begin{remark}[Perspectives within this volume]\label{rem:b2:metric:perspectives}
This chapter is the volume's load-bearing wall; watch where
each pillar carries weight. \emph{\omterm{def:b2:metric:complete}{Completeness}}: the Cauchy
criterion becomes the convergence test for series in Banach
spaces (\cref{ch:b2:series}), uniform convergence in
\cref{ch:b2:funcseq} is exactly convergence in the \omterm{def:b2:metric:complete}{complete}
$\bigl(C, d_\infty\bigr)$, and Cauchy--Lipschitz
(\cref{ch:b2:diffeq}) is the Banach fixed point theorem wearing
an integral equation. \emph{\omterm{def:b2:metric:compact}{Compactness}}: it proves the
equivalence of norms (\cref{ch:b2:nvs}), the attainment of
extrema for the optimization of \cref{ch:b2:diffcalc}, and the
existence of best approximations (\cref{ch:b2:nvs}'s weekend
problem). \emph{\omterm{def:b2:metric:connected}{Connectedness}}: it globalizes local statements
--- uniqueness of solutions of differential equations, the
intermediate value theorem on curves (\cref{ch:b2:curves}), and
the two components of $GL_n(\R)$ that orientation theory
(\cref{ch:b2:multint}) will keep apart.
\end{remark}

\begin{remark}[Common pitfalls]\label{rem:b2:metric:pitfalls}
(i) ``Closed and bounded implies \omterm{def:b2:metric:compact}{compact}'' is a theorem about
$\R^n$, not about \omterm{def:b2:metric:def}{metric spaces}: an infinite set with the
discrete metric is closed and bounded in itself yet not \omterm{def:b2:metric:compact}{compact}
(\cref{exo:b2:metric:4}), and the closed unit ball of
$C(\intcc01)$ fails too (\cref{ch:b2:nvs}). (ii) \omterm{def:b2:metric:complete}{Completeness} is
a property of the \emph{distance}, not of the topology: $\R$
with $d(x,y) = \abs{\arctan x - \arctan y}$ has the usual
convergent sequences but is incomplete
(\cref{exo:b2:metric:1}). (iii) A \omterm{def:b2:metric:continuity}{continuous} bijection need not
be a homeomorphism --- the circle parametrization of
\cref{exo:b2:metric:7}; \omterm{def:b2:metric:compact}{compactness} of the source repairs it.
(iv) Banach's theorem needs $k < 1$ \emph{uniformly}: the
condition $d(f(x), f(y)) < d(x,y)$ alone guarantees nothing on a
\omterm{def:b2:metric:compact}{non-compact} space (\cref{exo:b2:metric:5}). (v) \omterm{def:b2:metric:connected}{Connected} does
not imply \omterm{def:b2:metric:connected}{path-connected} in general --- but for the \omterm{def:b2:metric:topology}{open}
subsets of normed spaces met in this book, the two agree
(\cref{ch:b2:nvs}).
\end{remark}

\begin{remark}[Where this chapter is used]
Everywhere in the analysis half. \omterm{def:b2:metric:complete}{Completeness} of
$C(\intcc{a}{b})$ powers the convergence theorems of
\cref{ch:b2:funcseq} and the Cauchy--Lipschitz theory of
\cref{ch:b2:diffeq} (this chapter's weekend problem proves the
local Picard--Lindel\"of theorem already); \omterm{def:b2:metric:compact}{compactness} gives the
equivalence of norms in finite dimension (\cref{ch:b2:nvs}) and
the existence of extrema in \cref{ch:b2:diffcalc}; \omterm{def:b2:metric:connected}{connectedness}
underlies the intermediate value arguments of \cref{ch:b2:realfun}
and the uniqueness globalization for differential equations. In
the Year 3 volume, \omterm{def:b2:metric:compact}{compactness} in function spaces (the
Arzel\`a--Ascoli theorem) and the Baire category theorem
(\cref{exo:b2:metric:12} here) become tools of daily use.
\end{remark}

\begin{omfigure}
\begin{tikzpicture}[xscale=6.4, yscale=0.8]
  \draw[omDef, line width=3pt] (0,3) -- (1,3);
  \draw[omDef, line width=3pt] (0,2) -- (0.3333,2);
  \draw[omDef, line width=3pt] (0.6667,2) -- (1,2);
  \foreach \a in {0, 0.6667} {
    \draw[omDef, line width=3pt] (\a,1) -- (\a + 0.1111,1);
    \draw[omDef, line width=3pt]
      (\a + 0.2222,1) -- (\a + 0.3333,1);
  }
  \foreach \a in {0, 0.2222, 0.6667, 0.8889} {
    \draw[omThm, line width=3pt] (\a,0) -- (\a + 0.037,0);
    \draw[omThm, line width=3pt]
      (\a + 0.0741,0) -- (\a + 0.1111,0);
  }
  \node[left, font=\small] at (0,3) {$C_0$};
  \node[left, font=\small] at (0,2) {$C_1$};
  \node[left, font=\small] at (0,1) {$C_2$};
  \node[left, font=\small] at (0,0) {$C_3$};
\end{tikzpicture}

{\small The first stages of the Cantor set
(\cref{exo:b2:metric:8}): each level deletes the \omterm{def:b2:metric:topology}{open} middle
third of every segment. The intersection $C = \bigcap_n C_n$ is
\omterm{def:b2:metric:compact}{compact}, has empty interior and length zero, yet is \omterm{def:b2:structures:countable}{equipotent}
to $\R$ --- and it returns as a \emph{fixed point} of a
contraction on sets in this chapter's weekend problem (question
22).}
\end{omfigure}

\section{Connectedness}

\begin{definition}\label{def:b2:metric:connected}
$X$ is \emph{connected}\index{connected space} when it admits no
partition into two nonempty \omterm{def:b2:metric:topology}{open} subsets --- equivalently, when its
only subsets both \omterm{def:b2:metric:topology}{open} and closed are $\emptyset$ and $X$. $X$ is
\emph{path-connected} when any two points are joined by a \omterm{def:b2:metric:continuity}{continuous}
map $\gamma \colon \intcc{0}{1} \to X$.
\end{definition}

\begin{theorem}\label{thm:b2:metric:connectedness}
\begin{enumerate}
  \item The \omterm{def:b2:metric:connected}{connected} subsets of $\R$ are exactly the intervals.
  \item A \omterm{def:b2:metric:continuity}{continuous} image of a \omterm{def:b2:metric:connected}{connected space} is \omterm{def:b2:metric:connected}{connected} ---
        whence the general intermediate value theorem: a \omterm{def:b2:metric:continuity}{continuous}
        real function on a \omterm{def:b2:metric:connected}{connected space} takes every value between
        any two of its values.
  \item \omterm{def:b2:metric:connected}{Path-connected} $\Rightarrow$ \omterm{def:b2:metric:connected}{connected}. (The converse fails
        in general; it holds for \omterm{def:b2:metric:topology}{open} subsets of normed spaces,
        \cref{ch:b2:nvs}.)
\end{enumerate}
\end{theorem}

\begin{proof}
(1) A non-interval $A$ misses some $z$ between two of its points:
$A = (A \cap \intoo{-\infty}{z}) \cup (A \cap \intoo{z}{+\infty})$
splits it into two nonempty \omterm{def:b2:metric:topology}{open} (in $A$) pieces. Conversely, let
$I$ be an interval and $I = U \cup V$ a partition into nonempty
relatively \omterm{def:b2:metric:topology}{open} sets; pick $a \in U$, $b \in V$, say $a < b$, and
set $s = \sup\,(U \cap \intcc{a}{b})$, a point of $\intcc{a}{b}
\subseteq I$. If $s \in U$: then $s \neq b$, and relative openness
of $U$ puts a whole interval around $s$ (intersected with $I$)
inside $U$ --- so points of $U \cap \intcc{a}{b}$ exceed $s$,
contradicting the supremum. If $s \in V$: relative openness of $V$
puts an interval $\intoo{s - r}{s + r} \cap I$ inside $V$; but the
supremum is adherent to $U \cap \intcc{a}{b}$, which must meet that
interval --- contradiction with $U \cap V = \emptyset$. (This is
the Year 1 clopen argument for $\R$, run inside $I$.)

(2) If $f(X) = U' \cup V'$ splits into nonempty relatively \omterm{def:b2:metric:topology}{open}
sets, then $X = f^{-1}(U') \cup f^{-1}(V')$ splits $X$
(\cref{thm:b2:metric:globalcontinuity}). IVT: $f(X) \subseteq \R$
is \omterm{def:b2:metric:connected}{connected}, hence an interval by (1).

(3) Suppose $X = U \cup V$, both nonempty \omterm{def:b2:metric:topology}{open}, and join $a \in U$
to $b \in V$ by a path $\gamma$: then $\gamma^{-1}(U),
\gamma^{-1}(V)$ split $\intcc{0}{1}$, contradicting (1).
\end{proof}

\begin{example}\label{ex:b2:metric:glnr}
$GL_n(\R)$ is not \omterm{def:b2:metric:connected}{connected}: $\det$ is \omterm{def:b2:metric:continuity}{continuous} (a polynomial in
the entries) onto $\R^*$, which is not \omterm{def:b2:metric:connected}{connected}; the preimages of
$\R_+^*$ and $\R_-^*$ split $GL_n(\R)$. (Each piece is in fact
\omterm{def:b2:metric:connected}{path-connected} --- a pleasant exercise beyond our needs.) By
contrast $GL_n(\C)$ \emph{is} \omterm{def:b2:metric:connected}{path-connected}:
\cref{exo:b2:metric:10}.
\end{example}

\begin{example}[A fixed point from connectedness alone]\label{ex:b2:metric:ivtfixed}
Every \omterm{def:b2:metric:continuity}{continuous} $f \colon \intcc01 \to \intcc01$ has a fixed
point --- no contraction hypothesis, no iteration. Consider
$g(x) = f(x) - x$, \omterm{def:b2:metric:continuity}{continuous} on the \omterm{def:b2:metric:connected}{connected} $\intcc01$:
\[
g(0) = f(0) \geq 0,
\qquad
g(1) = f(1) - 1 \leq 0 ,
\]
and the intermediate value theorem
(\cref{thm:b2:metric:connectedness}~(2)) delivers a zero of
$g$, i.e.\ a fixed point of $f$. Contrast with Banach
(\cref{thm:b2:metric:banach}): here existence is topological
and free, but uniqueness and the algorithm are lost --- $f =
\mathrm{id}$ has every point fixed, and iteration of a
non-contracting $f$ may cycle forever. The two fixed point
theorems of this chapter answer different questions with
different currencies.
\end{example}

\begin{example}[$\R$ and $\R^2$ are not homeomorphic]\label{ex:b2:metric:rnotr2}
\omterm{def:b2:metric:connected}{Connectedness} is a topological fingerprint. Suppose $h \colon
\R^2 \to \R$ were a homeomorphism (a \omterm{def:b2:metric:continuity}{continuous} bijection with
\omterm{def:b2:metric:continuity}{continuous} inverse). Remove one point $a \in \R^2$: the
restriction $h \colon \R^2\setminus\{a\} \to
\R\setminus\{h(a)\}$ is still a homeomorphism. But $\R^2$ minus
a point is \omterm{def:b2:metric:connected}{path-connected} --- join any two points by a segment,
detouring along a second segment through an auxiliary point if
$a$ blocks the direct one --- hence \omterm{def:b2:metric:connected}{connected}
(\cref{thm:b2:metric:connectedness}~(3)); while $\R$ minus a
point splits into two nonempty \omterm{def:b2:metric:topology}{open} half-lines: not \omterm{def:b2:metric:connected}{connected}.
\omterm{def:b2:metric:connected}{Connectedness} is preserved by \omterm{def:b2:metric:continuity}{continuous} maps: contradiction.
The plane and the line are genuinely different \emph{as
topological spaces} --- a fact that cardinality alone
(\cref{exo:b2:structures:3}-style bijections do exist!) is too
coarse to see.
\end{example}

\section{Exercises}

\begin{exercise}[$\star$]\label{exo:b2:metric:1}
On $\R$, check that $\delta(x, y) = \min(1, \abs{x - y})$ and
$d(x,y) = \abs{\arctan x - \arctan y}$ are distances. Which
sequences converge for each? Is $(\R, d)$ \omterm{def:b2:metric:complete}{complete}?
\end{exercise}

\begin{exercise}[$\star$]\label{exo:b2:metric:2}
In a \omterm{def:b2:metric:def}{metric space}, prove that a convergent sequence is Cauchy and
bounded, and that a Cauchy sequence with a convergent subsequence
converges. Deduce again that \omterm{def:b2:metric:compact}{compact metric spaces} are \omterm{def:b2:metric:complete}{complete}.
\end{exercise}

\begin{exercise}[$\star$]\label{exo:b2:metric:3}
In $\bigl(C(\intcc{0}{1}), d_\infty\bigr)$, compute the distance
between $f(x) = x$ and $g(x) = x^2$; describe the closed ball
$\overline B(0, 1)$; and prove that the set $\{f : f(0) = 0\}$ is
closed while $\{f : f(0) > 0\}$ is \omterm{def:b2:metric:topology}{open}.
\end{exercise}

\begin{exercise}[$\star\star$]\label{exo:b2:metric:4}
Prove that the discrete \omterm{def:b2:metric:def}{metric space} $X$ (any set) is \omterm{def:b2:metric:complete}{complete}, and
that it is \omterm{def:b2:metric:compact}{compact} if and only if $X$ is finite. Which subsets are
\omterm{def:b2:metric:connected}{connected}?
\end{exercise}

\begin{exercise}[$\star\star$]\label{exo:b2:metric:5}
Let $X$ be \omterm{def:b2:metric:compact}{compact} and $f \colon X \to X$ with
\[
d\bigl(f(x), f(y)\bigr) < d(x, y) \quad \text{for all } x \neq y .
\]
Prove that $f$ has a unique fixed point \emph{(minimize $x \mapsto
d(x, f(x))$)}, and give an example on $X = \intco{1}{+\infty}$
(not \omterm{def:b2:metric:compact}{compact}) with no fixed point.
\end{exercise}

\begin{exercise}[$\star\star$]\label{exo:b2:metric:6}
(Nested \omterm{def:b2:metric:compact}{compacts}) Let $(K_n)$ be a decreasing sequence of nonempty
\omterm{def:b2:metric:compact}{compact} subsets of a \omterm{def:b2:metric:def}{metric space}. Prove that $\bigcap_n K_n \neq
\emptyset$ \emph{(pick $x_n \in K_n$ and extract)}. Show by example
that nonempty nested \emph{closed} sets in $\R$ can have empty
intersection.
\end{exercise}

\begin{exercise}[$\star\star$]\label{exo:b2:metric:7}
Let $K$ be \omterm{def:b2:metric:compact}{compact} and $f \colon K \to Y$ \omterm{def:b2:metric:continuity}{continuous} and bijective.
Prove that $f^{-1}$ is \omterm{def:b2:metric:continuity}{continuous} \emph{(use closed sets:
\cref{thm:b2:metric:globalcontinuity} and
\cref{thm:b2:metric:compactprops})}. Give a counterexample without
\omterm{def:b2:metric:compact}{compactness} ($\gamma(t) = (\cos t, \sin t)$ on
$\intco{0}{2\pi}$).
\end{exercise}

\begin{exercise}[$\star\star$]\label{exo:b2:metric:8}
The \emph{Cantor set} $C$ is obtained from $\intcc{0}{1}$ by
repeatedly deleting \omterm{def:b2:metric:topology}{open} middle thirds. Prove that $C$ is \omterm{def:b2:metric:compact}{compact},
has empty interior, and is infinite --- indeed \omterm{def:b2:structures:countable}{equipotent} to
$\{0,1\}^{\N}$ \emph{(ternary expansions with digits $0,2$;
\cref{exo:b2:structures:3})}.
\end{exercise}

\begin{exercise}[$\star\star\star$]\label{exo:b2:metric:9}
Let $X$ be a \emph{\omterm{def:b2:metric:compact}{compact}} \omterm{def:b2:metric:def}{metric space} and $f \colon X \to X$ an
isometry: $d(f(x), f(y)) = d(x, y)$. Prove that $f$ is surjective.
\emph{Hint: if $a \notin f(X)$, then $\varepsilon = d(a, f(X)) > 0$
(why?); study the orbit $a, f(a), f^2(a), \dots$ and show its
points are pairwise $\geq \varepsilon$ apart --- contradiction with
\omterm{def:b2:metric:compact}{compactness}.}
\end{exercise}

\begin{exercise}[$\star\star\star$]\label{exo:b2:metric:10}
Prove that $GL_n(\C)$ is \omterm{def:b2:metric:connected}{path-connected}. \emph{Hint: given $A, B$
invertible, consider $p(z) = \det\bigl((1 - z)A + zB\bigr)$ for $z
\in \C$: a polynomial in $z$, not identically zero, so it has
finitely many roots; join $0$ to $1$ in $\C$ by a path avoiding
them.}
\end{exercise}

\begin{exercise}[$\star\star$]\label{exo:b2:metric:11}
For $\emptyset \neq A \subseteq X$, set $d(x, A) = \inf_{a \in A}
d(x, a)$. Prove that $x \mapsto d(x, A)$ is $1$-Lipschitz, that
$d(x, A) = 0$ iff $x \in \overline A$, and that for disjoint
nonempty \emph{closed} sets $A, B$ the function
\[
\varphi(x) = \frac{d(x, A)}{d(x, A) + d(x, B)}
\]
is well defined, \omterm{def:b2:metric:continuity}{continuous}, equal to $0$ exactly on $A$ and to
$1$ exactly on $B$ --- a \omterm{def:b2:metric:continuity}{continuous} ``switch'' separating any two
disjoint closed sets.
\end{exercise}

\begin{exercise}[$\star\star\star$]\label{exo:b2:metric:12}
(Baire) Let $X$ be a \omterm{def:b2:metric:complete}{complete metric space} and $(U_n)_{n\geq1}$ a
sequence of dense \omterm{def:b2:metric:topology}{open} subsets. Prove that $\bigcap_n U_n$ is
dense in $X$ \emph{(inside any ball, build nested closed balls
$\overline B(x_n, r_n) \subseteq U_n$ with $r_n \to 0$ and use
\omterm{def:b2:metric:complete}{completeness})}. Deduce that $\R$ is not a \omterm{def:b2:structures:countable}{countable} union of
closed sets with empty interior, and recover --- again --- that
$\R$ is uncountable.
\end{exercise}

\section{Problem: Picard Iteration}

\omterm{def:b2:metric:complete}{Completeness} plus contraction is a solving machine: feed it an
equation written as a fixed point problem, and it returns
existence, uniqueness, an algorithm, and error bars. This weekend
problem runs the machine at full power on the equation $y' = f(t,
y)$: we prove the local \emph{Picard--Lindel\"of theorem} (the
nonlinear heart of \cref{ch:b2:diffeq}'s Cauchy--Lipschitz
theory), watch every hypothesis earn its keep through
counterexamples, and collect purely metric dividends ---
\omterm{def:b2:metric:continuity}{continuous} dependence on the data, Kepler's equation, and the
self-similarity of the Cantor set.

\begin{omfigure}
\begin{tikzpicture}
\begin{axis}[omaxis, width=8.6cm, height=5.4cm,
    xmin=-0.05, xmax=1.45, ymin=0, ymax=4.2,
    xtick={0,1}, ytick={1,2,3}]
  \addplot[omDef!45, thick, domain=0:1.4, samples=2] {1};
  \addplot[omDef!65, thick, domain=0:1.4, samples=2] {1 + x};
  \addplot[omDef, thick, domain=0:1.4, samples=60]
    {1 + x + x^2/2};
  \addplot[omThm, very thick, domain=0:1.4, samples=120]
    {exp(x)};
  \node[font=\small, omDef] at (axis cs:1.28,0.8) {$y_0$};
  \node[font=\small, omDef] at (axis cs:1.28,2.15) {$y_1$};
  \node[font=\small, omDef] at (axis cs:1.28,2.95) {$y_2$};
  \node[font=\small, omThm] at (axis cs:1.17,3.85)
    {$\eu^{\,t}$};
\end{axis}
\end{tikzpicture}

{\small Picard iterates for $y' = y$, $y(0) = 1$: each pass
through $T(y)(t) = 1 + \int_0^t y$ adds one Taylor term, and the
contraction squeezes the whole sequence uniformly onto
$\eu^{\,t}$.}
\end{omfigure}

\begin{problem}[{Weekend problem --- the Picard--Lindel\"of
theorem}]
\label{pb:b2:metric:1}
Throughout, $t_0 \in \R$, $y_0 \in \R$, $a, b > 0$, and $f$ is a
\omterm{def:b2:metric:continuity}{continuous} function on the rectangle $R = \intcc{t_0 - a}{t_0 +
a} \times \intcc{y_0 - b}{y_0 + b}$, bounded by $M =
\sup_R\,\abs f$, and \emph{$L$-Lipschitz in its second variable}:
$\abs{f(t, y) - f(t, z)} \leq L\abs{y - z}$ whenever both points
lie in $R$. Set
\[
h = \min\Bigl(a, \frac bM\Bigr) \quad (\text{with } h = a
\text{ if } M = 0), \qquad I = \intcc{t_0 - h}{t_0 + h}.
\]

\medskip
\noindent\textbf{Part I --- The \omterm{def:b2:metric:complete}{complete} stage.}
\begin{enumerate}
  \item Prove the two statements quoted in
        \cref{def:b2:metric:complete}: a closed subset of a
        \omterm{def:b2:metric:complete}{complete metric space} is \omterm{def:b2:metric:complete}{complete}, and a \omterm{def:b2:metric:complete}{complete} subset
        of any \omterm{def:b2:metric:def}{metric space} is closed. Deduce that every closed
        subset of $\bigl(C(I), d_\infty\bigr)$ is a \omterm{def:b2:metric:complete}{complete
        metric space}.
  \item Show that the map $C(I) \to C(I)$, $y \mapsto \bigl(t
        \mapsto \int_{t_0}^{t}y(s)\,\dd s\bigr)$, is
        $h$-Lipschitz for $d_\infty$.
  \item (Fixed points move less than the maps) Let $g \colon X
        \to X$ be a $k$-contraction of a \omterm{def:b2:metric:def}{metric space} with fixed
        point $\ell_g$, and $\widetilde g \colon X \to X$
        \emph{any} map with a fixed point $\ell_{\widetilde g}$.
        Prove
        \[
        d(\ell_g, \ell_{\widetilde g}) \leq
        \frac{d\bigl(g(\ell_{\widetilde g}),
        \widetilde g(\ell_{\widetilde g})\bigr)}{1 - k}
        \leq \frac{\sup_{x \in X} d\bigl(g(x), \widetilde
        g(x)\bigr)}{1 - k}.
        \]
  \item (The iterate trick) Let $X$ be \omterm{def:b2:metric:complete}{complete} nonempty and $g
        \colon X \to X$ a map --- not assumed \omterm{def:b2:metric:continuity}{continuous} ---
        such that some iterate $g^m$ is a $k$-contraction.
        Prove that $g$ has a unique fixed point $\ell$ and that
        \emph{every} orbit $x_{n+1} = g(x_n)$ converges to
        $\ell$. \emph{(Fixed points of $g$ are fixed points of
        $g^m$; conversely $g(\ell)$ is a fixed point of $g^m$;
        split the orbit along residues mod $m$.)}
\end{enumerate}

\medskip
\noindent\textbf{Part II --- The Picard--Lindel\"of theorem.}
\begin{enumerate}[resume]
  \item Show that a function $y \colon I \to \intcc{y_0 - b}{y_0
        + b}$ is $C^1$ with $y(t_0) = y_0$ and $y' = f(t, y)$ on
        $I$ if and only if it is \omterm{def:b2:metric:continuity}{continuous} and satisfies the
        integral equation
        \[
        y(t) = y_0 + \int_{t_0}^{t} f\bigl(s, y(s)\bigr)\dd s
        \qquad (t \in I).
        \]
  \item Let $X_h = \{y \in C(I) : \abs{y(t) - y_0} \leq b
        \text{ on } I\}$ and let $T$ be defined by $T(y)(t) =
        y_0 + \int_{t_0}^{t}f(s, y(s))\dd s$. Show that $X_h$ is
        a nonempty closed subset of $C(I)$, hence \omterm{def:b2:metric:complete}{complete}, and
        that $T$ maps $X_h$ into $X_h$ --- this is where $h \leq
        b/M$ works.
  \item Show that $d_\infty\bigl(T(y), T(z)\bigr) \leq
        Lh\,d_\infty(y, z)$ on $X_h$: if $Lh < 1$, Banach's
        theorem already concludes. We remove this smallness
        condition next.
  \item Prove by induction on $n$:
        \[
        \abs{T^n(y)(t) - T^n(z)(t)} \leq
        \frac{\bigl(L\abs{t - t_0}\bigr)^n}{n!}\,
        d_\infty(y, z) \qquad (y, z \in X_h,\ t \in I),
        \]
        so some iterate of $T$ is a contraction. Conclude with
        question 4 (\emph{the Picard--Lindel\"of theorem}): the
        Cauchy problem $y' = f(t,y)$, $y(t_0) = y_0$ has exactly
        one solution on $I = \intcc{t_0 - h}{t_0 + h}$ with
        values in $\intcc{y_0 - b}{y_0 + b}$.
  \item Show that the restriction ``with values in $\intcc{y_0 -
        b}{y_0 + b}$'' is automatic: any solution of the Cauchy
        problem defined on $I$ whose graph starts in $R$ stays
        in $\intcc{y_0 - b}{y_0 + b}$ \emph{(consider the first
        exit time and bound $\abs{y(t) - y_0}$ by $M\abs{t -
        t_0}$)}. Hence uniqueness holds among all solutions on
        $I$.
  \item Run the machine on $y' = y$, $y(0) = 1$, starting from
        the constant $y^{(0)} \equiv 1$: compute the Picard
        iterates $y^{(n)}$, identify them, and describe the
        convergence.
\end{enumerate}

\medskip
\noindent\textbf{Part III --- Every hypothesis earns its keep.}
\begin{enumerate}[resume]
  \item (\omterm{def:b2:metric:continuity}{Lipschitz} fails, uniqueness fails) For $y' =
        2\sqrt{\abs y}$, $y(0) = 0$: check that $y \equiv 0$
        and, for every $c \geq 0$, the function $y_c(t) = 0$ for
        $t \leq c$, $y_c(t) = (t - c)^2$ for $t > c$, are all
        $C^1$ solutions on $\R$. Where exactly does $y \mapsto
        2\sqrt{\abs y}$ fail to be \omterm{def:b2:metric:continuity}{Lipschitz}?
  \item (Locality is real) For $y' = y^2$, $y(0) = 1$: solve
        explicitly, give the maximal interval of existence, and
        compute the best $h$ the theorem can certify over all
        choices of the rectangle ($a$ large, $b$ free): show
        $h_{\max} = \sup_{b>0} \frac{b}{(1+b)^2} = \frac14$,
        while the true solution lives on $\intoo{-\infty}{1}$.
  \item (\omterm{def:b2:metric:complete}{Completeness} is not decor) On $X = \Q \cap
        \intcc{1}{2}$ with the usual distance, let $g(x) =
        \frac x2 + \frac1x$. Show $g(X) \subseteq X$, that $g$
        is a $\frac12$-contraction \emph{(mean value
        inequality)}, and that $g$ has no fixed point in $X$.
        Which hypothesis of Banach's theorem fails, and what is
        the fixed point in the completion?
  \item (Error bars) For a $k$-contraction $g$ on a \omterm{def:b2:metric:complete}{complete}
        space, prove the \emph{a posteriori} estimate $d(x_n,
        \ell) \leq \frac{k}{1-k}\,d(x_n, x_{n-1})$. For Heron's
        map $g(x) = \frac x2 + \frac 1x$ on $\intcc{1}{2}$
        (fixed point $\sqrt2$), starting at $x_0 = \frac32$:
        how many steps does the \emph{a priori} bound
        $\frac{k^n}{1-k}d(x_1, x_0)$ demand for accuracy
        $10^{-6}$, and how many steps suffice in reality?
        (Compute $x_1, x_2, x_3$ and their errors; the
        contraction bound is honest but pessimistic ---
        Heron converges quadratically.)
\end{enumerate}

\medskip
\noindent\textbf{Part IV --- \omterm{def:b2:metric:continuity}{Continuous} dependence.} In this
part $Lh < 1$, so $T$ itself is a contraction on $X_h$ (question
7); write $y[\,y_0\,]$ for the solution with initial value
$y_0$.
\begin{enumerate}[resume]
  \item (Dependence on the initial value) Let $z_0$ be another
        initial value with $\abs{z_0 - y_0}$ small enough that
        both problems fit in the rectangle. Using question 3,
        prove
        \[
        d_\infty\bigl(y[y_0], y[z_0]\bigr) \leq
        \frac{\abs{y_0 - z_0}}{1 - Lh}.
        \]
  \item (Long intervals by chaining) Suppose the solutions exist
        on a long segment cut into $m$ consecutive pieces on
        each of which the previous bound applies with $Lh \leq
        \frac12$. Show the deviation grows by a factor at most
        $2$ per piece, hence $d_\infty \leq 2^m\abs{y_0 - z_0}$
        overall --- an exponential-in-length bound, the discrete
        shadow of the $\eu^{L\abs{t - t_0}}$ of Gronwall's lemma
        (\cref{ch:b2:diffeq}).
  \item (Dependence on the field) Let $g$ be another field on
        $R$, also $L$-Lipschitz in $y$, with $\sup_R \abs{f - g}
        \leq \varepsilon$. Prove that the corresponding
        solutions satisfy $d_\infty \leq
        \frac{\varepsilon h}{1 - Lh}$: modelling error
        propagates linearly.
  \item (Parameters) If a family $f_\lambda$ of fields is
        $L$-Lipschitz in $y$ uniformly and $\sup_R\abs{f_\lambda
        - f_\mu} \leq C\abs{\lambda - \mu}$, deduce that
        $\lambda \mapsto y_\lambda$ is \omterm{def:b2:metric:continuity}{Lipschitz} from the
        parameter space into $\bigl(C(I), d_\infty\bigr)$.
  \item (Systems cost nothing) Explain why Parts I, II and IV
        hold verbatim for $y$ with values in $\R^n$ (sup
        distances built on any of the distances of
        \cref{ex:b2:metric:examples}), then compute all Picard
        iterates for the system $y' = Ay$, $y(0) = (c_1,
        c_2)$, $A = \begin{psmallmatrix} 0 & 1\\ 0 &
        0\end{psmallmatrix}$: show the iteration becomes
        stationary at the exact solution after one step.
\end{enumerate}

\medskip
\noindent\textbf{Part V --- Metric dividends and synthesis.}
\begin{enumerate}[resume]
  \item (Perturbation of the identity) Let $X$ be a \omterm{def:b2:metric:complete}{complete}
        normed-space-like stage: take $X = C(I)$ or $\R^n$. If
        $\eta \colon X \to X$ is $k$-Lipschitz with $k < 1$,
        prove that $x \mapsto x + \eta(x)$ is a bijection of $X$
        whose inverse is $\frac1{1-k}$-Lipschitz \emph{(for each
        $y$, apply Banach to $x \mapsto y - \eta(x)$)}. This is
        the metric heart of the inverse function theorem
        (\cref{ch:b2:diffcalc}).
  \item (Kepler's equation) For $0 \leq e < 1$ and $m \in \R$,
        prove that $x = m + e\sin x$ has exactly one solution,
        that the iteration $x_{n+1} = m + e\sin x_n$ converges
        to it from any start, and estimate: for $e = \frac12$,
        $m = 1$, how many iterations guarantee an error $\leq
        10^{-3}$ by the a priori bound? (The solution is $x
        \approx 1.4987$.)
  \item (The Cantor set is a fixed point) Let $S_1(x) = \frac
        x3$ and $S_2(x) = \frac x3 + \frac23$ on $\R$, and let
        $C$ be the Cantor set of \cref{exo:b2:metric:8}. Prove
        $C = S_1(C) \cup S_2(C)$, and explain in one sentence
        why no \emph{other} nonempty \omterm{def:b2:metric:compact}{compact} set satisfies this
        equation (the map $A \mapsto S_1(A) \cup S_2(A)$ is a
        contraction for a distance between \omterm{def:b2:metric:compact}{compact} sets --- the
        Hausdorff distance, made honest in the Year 3 volume).
  \item (\omterm{def:b2:metric:connected}{Connectedness} globalizes uniqueness) Let $f$ be locally
        \omterm{def:b2:metric:continuity}{Lipschitz} in $y$ on an \omterm{def:b2:metric:topology}{open} set, and let $y, z$ be two
        solutions of $y' = f(t, y)$ on a common interval $J$
        with $y(t_0) = z(t_0)$. Prove $y = z$ on $J$: show that
        $\{t \in J : y(t) = z(t)\}$ is nonempty, closed in $J$,
        and \omterm{def:b2:metric:topology}{open} in $J$ (by local uniqueness), and use the
        \omterm{def:b2:metric:connected}{connectedness} of intervals
        (\cref{thm:b2:metric:connectedness}).
  \item (No smallness for linear equations) For $y' =
        \alpha(t)y + \beta(t)$ with $\alpha, \beta$ \omterm{def:b2:metric:continuity}{continuous}
        on a segment $\intcc{A}{B}$, adapt question 8 to show
        the factorial bound holds on the \emph{whole} segment,
        so existence and uniqueness are global there --- the
        scalar case of \cref{ch:b2:diffeq}'s Cauchy--Lipschitz
        theorem, with no restriction on the length $B - A$.
  \item (Synthesis) One sentence each: what \omterm{def:b2:metric:complete}{completeness}
        contributed; what the contraction contributed; what the
        iterate trick bought compared with plain Banach; where
        \omterm{def:b2:metric:connected}{connectedness} entered; and which counterexample of Part
        III guards which hypothesis. Name the summit theorem,
        and say what replaces contraction when $f$ is merely
        \omterm{def:b2:metric:continuity}{continuous} (Peano's theorem, via \omterm{def:b2:metric:compact}{compactness} in function
        spaces --- the Year 3 volume's Arzel\`a--Ascoli).
\end{enumerate}
\end{problem}
